\chapter{Az eredmények értékelése}
% 5.1
%~~~~~~~~~~~~~~~~~~~~~~~~~~~~~~~~~~~~~~~~~~~~~~~~~~~~~~~~~~~~~~~~~~~~~~~~~~~~~~~~~~~~~~
\section{Rendszer- és teljesítményelemzés (Kritikai reflexió)}
	Az eseményvezérelt architektúra hatása

	Az InfluxDB HTTP kapcsolat overheadje?

% 5.2
%~~~~~~~~~~~~~~~~~~~~~~~~~~~~~~~~~~~~~~~~~~~~~~~~~~~~~~~~~~~~~~~~~~~~~~~~~~~~~~~~~~~~~~
\section{A döntési modellek kísérleti összehasonlítása}
	Melyik gondolkodásmód bizonyult a legstabilabbnak?

	A genetikai drift elemzése: hogyan változtak a populáció átlagértékei X generáció után

% 5.3
%~~~~~~~~~~~~~~~~~~~~~~~~~~~~~~~~~~~~~~~~~~~~~~~~~~~~~~~~~~~~~~~~~~~~~~~~~~~~~~~~~~~~~~
\section{Továbbfejlesztési lehetőségek}

A kifejlesztett ökoszisztéma-szimuláció egy stabil, moduláris alapkeretet biztosít, amely mind szoftverarchitekturális, mind biológiai modellezési szempontból jelentős továbbfejlesztési potenciállal rendelkezik. A jövőbeli kutatási és fejlesztési irányok három fő kategóriába csoportosíthatók.

\subsection{Technológiai irány: Skálázhatóság és a Unity DOTS architektúra}

A jelenlegi implementáció a hagyományos objektumorientált (\textit{Object-Oriented Programming, OOP}) megközelítést alkalmazza, ahol minden ágens egy egyedi \texttt{MonoBehaviour} alapú \texttt{GameObject} a Unity színtéren. Bár ez a struktúra ideális a komplex állapotgépek és a Fuzzy Cognitive Map egyedi kiértékelésére, komoly teljesítménybeli korlátokat (\textit{performance bottleneck}) jelent a CPU és a memória-hozzáférés szintjén. A szimulált ágensek száma a hardverkapacitástól függően néhány száz, maximum egy-kétezer egyedben maximalizálódik, mivel az objektumok kezelése és a Unity belső alrendszereinek (pl. \texttt{Update()} ciklusok overheadje) jelentős erőforrást igényel.

A szimuláció nagyságrendi skálázhatóságának biztosításához a technológiai jövőkép a \textbf{Unity DOTS (Data-Oriented Technology Stack)} és az \textbf{ECS (Entity Component System)} architektúra bevezetése.

\begin{itemize}
    \item \textbf{Adatközpontú tervezés (ECS):} A \texttt{GameObject}-ek helyett az egyedeket egyszerű azonosítók (\textit{Entities}), a tulajdonságaikat (pozíció, éhség, genetikai súlyok) tiszta adatstruktúrák (\textit{Components}), a logikát pedig globális rendszerek (\textit{Systems}) reprezentálnák. Ez megszüntetné a memóriatöredezettséget, és lehetővé tenné a gyorsítótár-optimalizált (\textit{cache-friendly}) szekvenciális adatolvasást.
    \item \textbf{Párhuzamosítás (C\# Job System \& Burst Compiler):} Az ECS-re való átállással az egyedek döntéshozatali ciklusai (legyen az hasznossági alapú vagy FCM hálózat) natív gépi kódra fordulnának, és automatikusan szétoszthatók lennének a CPU összes elérhető magja között.
\end{itemize}

Ezen technológiai paradigmaváltás eredményeképpen a szimulációs kapacitás a jelenlegi néhány száz ágensről \textbf{több tízezer párhuzamosan működő ágensre} növekedne, ami már valódi, makroszintű populációdinamikai és makroevolúciós trendek vizsgálatát is lehetővé tenné.

\subsection{Biológiai és ökológiai irány: A szimulációs hűség növelése}

A környezeti és viselkedési modell finomhangolásával a szimuláció közelebb vihető a valós ökoszisztémák komplexitásához. A legfontosabb biológiai fejlesztési irányok a következők:

\paragraph{Dinamikus környezeti tényezők:}
\begin{itemize}
    \item \textbf{Nappal-éjszaka ciklusok:} Az ágensek látási sugarának (\texttt{sensor.radius}) és camouflage mutatóinak időfüggő skálázása. Ez lehetőséget adna az éjszakai (nokturnális) és nappali (diurnális) életmód evolúciós kialakulására.
    \item \textbf{Szezonalitás (Évszakok):} A növényzet reprodukciós rátájának és a vízforrások elérhetőségének globális, ciklikus változtatása, amely rákényszerítené a populációt a vándorlásra, az erőforrások felhalmozására vagy a reprodukciós időszakok optimalizálására.
    \item \textbf{Domborzatfüggő mikroklímák:} Az \textit{Island} terep magassági és koordináta-adatai alapján különböző szárazsági vagy hőmérsékleti zónák kialakítása, amelyek befolyásolnák a staminaveszteséget vagy a metabolikus rátát.
\end{itemize}

\paragraph{Komplex társas és szaporodási interakciók:}
\begin{itemize}
    \item \textbf{Falka-viselkedés (Flocking/Swarming):} Csoportos vadászati vagy védekezési stratégiák kódolása. A ragadozók (rókák) képesek lennének koordináltan bekeríteni a zsákmányt, míg a prédák (nyulak) csoportba tömörülve csökkenthetnék az egyéni predációs kockázatot.
    \item \textbf{Utódgondozás és fészekrakás:} A szaporodási mechanizmus kiterjesztése fix vagy dinamikusan épített objektumokkal (pl. nyúlüregek, kotorékok). A vemhes nőstények nem a nyílt színtéren hoznák világra az utódokat, hanem védett környezetben, ahol a kölykök egy bizonyos fejlődési fázisig (\texttt{currentAge < 1}) védve lennének a predációtól, miközben a szülők táplálékot hordanak számukra.
\end{itemize}

Ezen irányok bevezetése nemcsak a szimuláció vizuális és funkcionális komplexitását növelné, hanem új dimenziókat nyitna az adaptív ágensek genetikai konvergenciájának és a mesterséges élet (\textit{Artificial Life}) kutatásának területén.

\subsection{Játékmenet és gamifikáció: Interaktív ökoszisztéma-szimulátor és Sandbox mód}

Mivel a keretrendszer natívan a Unity játékmotorra épül, a szimuláció funkcionális kiterjesztésének egy harmadik, magától értetődő iránya a projekt gamifikációja (\textit{gamification}). Ez a megközelítés a tiszta tudományos modellezést egy interaktív, oktatási vagy szórakoztató célú „God game” vagy túlélőjáték-mechanizmussá alakítaná át.

\paragraph{Interaktív Sandbox és paraméter-tuning:}
A jövőbeli fejlesztések során megvalósítható egy olyan felhasználói felület (UI), amelyen keresztül a felhasználó valós időben, közvetlenül beavatkozhat az ökoszisztéma működésébe. A játékos maga konfigurálhatná a kezdeti génkészleteket, manuálisan helyezhetne el új egyedeket, vagy akár globális környezeti katasztrófákat (pl. aszály, járványok) idézhetne elő, tesztelve a populációk túlélőképességét és evolúciós stabilitását.

\paragraph{Közvetlen ágens-irányítás és túlélési mód:}
A gamifikáció legmagasabb szintje egy olyan hibrid játékmód bevezetése lenne, ahol a felhasználó átvehetné az irányítást egy kiválasztott ágens (például egy nyúl vagy egy róka) felett. Ebben a túlélési módban a játékosnak magának kellene menedzselnie az egyed belső szükségleteit (éhség, szomjúság, staminaveszteség), miközben a környezetében lévő többi ágens továbbra is autonóm módon működik.

\paragraph{Döntéshozatali architektúrák aszimmetrikus versenye:}
Ez a megközelítés egy egyedülálló, aszimmetrikus versenyt eredményezne a humán intelligencia és a beépített mesterséges intelligencia-modellek között. A játékos közvetlenül összemérhetné túlélési stratégiáját és reflexeit:
\begin{itemize}
    \item a merev, szabályalapú, de kiszámítható \texttt{ClassicFSM} ágensekkel,
    \item a pillanatnyi prioritásokat matematikailag súlyozó \texttt{Utility}-alapú egyedekkel, valamint
    \item a komplex, asszociatív visszacsatolásokat és félelmet is modellezni képes, evolválódott \texttt{FCM} hálózatokkal.
\end{itemize}

Ez a fejlesztési irány nemcsak a szoftver demonstrációs értékét növelné, hanem kiváló alapként szolgálhatna oktatási szoftverekhez (pl. az ökológiai egyensúly és a természetes szelekció szemléltetésére), vagy akár egy komplex, önálló stratégiai játék prototípusaként is megállná a helyét.

% 5.4
%~~~~~~~~~~~~~~~~~~~~~~~~~~~~~~~~~~~~~~~~~~~~~~~~~~~~~~~~~~~~~~~~~~~~~~~~~~~~~~~~~~~~~~
\section{Összefoglalás}

Jelen diplomamunka egy négy féléven átívelő, szisztematikus kutatási és szoftverfejlesztési folyamat szintézise, melynek elsődleges célja egy olyan robusztus és moduláris ökoszisztéma-szimulációs keretrendszer létrehozása volt, amely alkalmas populációgenetikai, evolúciós és ágensalapú döntéshozatali modellek tudományos igényű vizsgálatára, valamint valós idejű telemetriai elemzésére.

A fejlesztési folyamat kiindulópontját képező BSc szakdolgozati fázisban egy korlátozott funkcionalitású, monolitikus struktúrájú alkalmazás állt rendelkezésre. Ez a korai rendszer kizárólag egyetlen préda faj autonóm működését modellezte kódgenerált környezetben, merev, determinisztikus állapotátmenetekre támaszkodva, lokális és strukturálatlan adatmentés mellett.

A mesterképzés (MSc) során elvégzett mérnöki munka alapvetően változtatta meg a szoftver irányvonalát: a korábbi korlátokat felszámolva a hangsúly a tiszta szoftverarchitekturális mintákra, a skálázhatóságra, a biológiai hűség növelésére és az elosztott adatgyűjtésre helyeződött át.

A szoftvertechnológiai modernizáció első lépéseként a „god class” antipatternek felszámolásával egy tiszta, komponens-aggregációra épülő osztályhierarchiát alakítottam ki a Unity3D motor környezetében. A számítási kapacitás optimalizálása és az ágensek közötti decoupling érdekében a rendszert eseményvezérelt kommunikációs modellre ültettem át. Ez a megközelítés minimalizálta a felesleges, frame-enként lefutó állapotellenőrzéseket, biztosítva a CPU-erőforrások hatékony elosztását növekvő egyedszám mellett is.

A biológiai hűség fokozása érdekében a szimulációs teret procedurálisan generált, folytonos 3D-s domborzattá terjesztettem ki, amely fizikai akadályaival közvetlen szelekciós nyomást gyakorol az ágensekre. Bevezetésre került a ragadozó faj (rókák), ezáltal sikerült digitális környezetben reprodukálni a klasszikus Lotka–Volterra-féle ragadozó-préda populáció-oszcillációkat. Az egyedek viselkedését egy komplex, mutációra és rekombinációra képes genetikai adatmodell (\texttt{Genome}) vezérli. A szuper-ágensek kialakulását versengő trade-off mechanizmusokkal akadályoztam meg, ahol a magasabb sebesség vagy látótávolság arányosan növeli az egyed metabolikus energiaköltségét.

A dolgozat egyik legfontosabb tudományos és mérnöki eredménye a döntéshozatali folyamatok moduláris absztrakciója, amely lehetővé tette három radikálisan eltérő döntési paradigma párhuzamos integrációját és összehasonlítását:
\begin{itemize}
    \item \textbf{Véges állapotgép (ClassicFSM):} A szabályalapú, stabil, de futásidőben merev és evolúciós szempontból alkalmatlan megközelítés reprezentációja.
    \item \textbf{Hasznossági alapú döntési rendszer (Utility AI):} Ahol az ágensek a normalizált belső szükségletek és külső ingerek lineáris kombinációjaként számítják ki az akciók prioritását. A súlyokat a \texttt{UtilityGenome} hordozza, folytonos adaptációs teret biztosítva a természetes szelekciónak.
    \item \textbf{Fuzzy Kognitív Térkép (FCM):} A legkomplexebb integrált modell, amely egy irányított, súlyozott gráf segítségével, iteratív és nem-lineáris szigmoid aktivációs függvényeken keresztül modellezi a belső állapotok (pl. éhség, félelem) és akciók egymásra hatását a \texttt{FCMGenome} súlymátrixa alapján.
\end{itemize}

A szimulációk során keletkező adatok tudományos kiértékeléséhez egy professzionális ipari telemetriai pipeline-t terveztem és valósítottam meg. A Unity főszálának védelme érdekében aszinkron architektúrát alkalmazva, HTTP POST kéréseken keresztül transzportálom a strukturált populációdinamikai és genetikai adatokat egy külső InfluxDB idősor-adatbázisba. Az adatok vizuális reprezentációját, valós idejű statisztikai elemzését és a különböző döntési modellek stabilitásának összehasonlítását egy egyedileg konfigurált Grafana dashboard rendszer biztosítja.

Összességében a kitűzött feladatokat hiánytalanul teljesítettem. A létrehozott szoftverrendszer nem csupán igazolta a biológiai alapelméletek digitális reprezentálhatóságát, hanem egy olyan skálázható és kiterjeszthető mérnöki keretrendszert biztosít, amely szilárd alapot nyújt a jövőbeli kutatásokhoz, legyen szó akár a Unity DOTS alapú masszív párhuzamosítással történő skálázásról, vagy a Sandbox módú gamifikációs továbbfejlesztésekről.
